% Solutions of Book 4, chapter 4 (cpp-templates).
\section*{Chapter \ref{ch:b2:cpp-templates} --- Templates and Generic Programming}

\begin{solution}{exo:b2:cpp-templates:1}
\omoutput{code/bachelor-2/04-cpp-templates/outputs/exo_overload.txt}{The seven lines.}
a: \code{show(int)} is an exact match and not a \omterm{def:b2:cpp-templates:template}{template}, so it wins over the \omterm{def:b2:cpp-templates:template}{template} with \code{T = int}. b: the
same with \code{show(double)}. c: \code{1.5f} is a \code{float}; \code{show(double)} needs a promotion, the \omterm{def:b2:cpp-templates:template}{template}
with \code{T = float} is exact. d: \code{'x'} is a \code{char}: exact for the \omterm{def:b2:cpp-templates:template}{template} with \code{T = char}, a
promotion for \code{show(int)}. e: explicit \omterm{def:b2:cpp-templates:template}{template} arguments leave only \omterm{def:b2:cpp-templates:template}{templates}, \code{T = int}. f: both
\omterm{def:b2:cpp-templates:template}{templates} match exactly (\code{T = int *} and \code{T = int}); the one with parameter \code{T *} is more specialised
and wins. g: \code{nullptr} has type \code{std::nullptr_t}, which converts to no \code{int} or \code{double} and from
which \code{T *} cannot be deduced; only \code{show<T>(T)} matches.
\end{solution}

\begin{solution}{exo:b2:cpp-templates:2}
\code{main_twice.cpp} sees only the declaration, so the compiler emits a call to \code{twice<int>} and trusts the
linker to find it. \code{twice.cpp} contains the \omterm{def:b2:cpp-templates:template}{template} but never uses it, so nothing there is instantiated, and
no file defines \code{twice<int>}: the linker reports it undefined. Fix 1, the usual one: define the \omterm{def:b2:cpp-templates:template}{template} in the
header. Fix 2: keep the definition in \code{twice.cpp} and add an \emph{explicit instantiation},
\code{template int twice<int>(int);}, for every argument list the program will use, which only works when that list
is known in advance.
\end{solution}

\begin{solution}{exo:b2:cpp-templates:3}
\omcode{code/bachelor-2/04-cpp-templates/exercises/clamp.hpp}{6}{19}{\code{clamp_to}}
It returns a reference to one of its arguments to avoid a copy of a large value. That is safe when the result is
used at once or copied, as in \code{int y = clamp_to(x, 0, 10);}. It dangles in \code{const int &r = clamp_to(x, 0,
10);} when the result is the temporary \code{0} or \code{10}, which dies at the end of the statement: the same caveat
as for the standard \code{std::clamp}. \code{tests/exercises_test.cpp} checks a type without \code{<} is refused,
through a \omterm{def:b2:cpp-templates:concept}{concept} that asks whether the call is valid.
\end{solution}

\begin{solution}{exo:b2:cpp-templates:4}
For each of \code{int}, \code{double} and \code{std::string}: the \omterm{def:b2:cpp-classes-raii:constructor}{constructor} and \omterm{def:b2:cpp-classes-raii:constructor}{destructor}, \code{push_back},
\code{reserve} and \code{clear}; plus \code{at} for \code{Vec<int>} and the \omterm{def:b2:cpp-copy-move:copy}{copy constructor} for
\code{Vec<std::string>}, and the \code{const} \code{operator[]} that this \omterm{def:b2:cpp-copy-move:copy}{copy constructor} calls on its
argument: 18 \omterm{def:b2:cpp-classes-raii:class}{member functions} (the defaulted \omterm{def:b2:cpp-classes-raii:constructor}{constructor} may generate no code). Members never used,
such as \code{capacity}, are not instantiated. Split in two files that both use \code{Vec<int>}, each object file
contains its own copies of the \code{Vec<int>} members, as weak symbols, and the linker keeps one of each: the
program contains the same 18 functions.
\end{solution}

\begin{solution}{exo:b2:cpp-templates:5}
1: \code{T} is deduced as \code{int} from 3 and as \code{double} from 7.5, a conflict; write \code{larger<double>(3,
7.5)} or \code{larger(3.0, 7.5)}. 2: \code{T} appears only in the return type, so nothing deduces it; write
\code{parse<int>("12")}. 3: the two arguments are arrays of 5 and 9 characters, which, passed by value, decay to
pointers: \code{T} is \code{const char *} for both, and the call compiles. It compares two \emph{addresses} of
unrelated string literals, whose order is unspecified, and returns one of the two pointers, converted to a string.
Fix: \code{larger(std::string("pier"), std::string("dunes rd"))}.
\end{solution}

\begin{solution}{exo:b2:cpp-templates:6}
\code{Ordered} asks only for \code{<}, and both types have it, so the call is accepted; the body also computes
\code{(lo + hi) / 2}, and for strings \code{+} concatenates but nothing divides a string by 2. Steps 1 and 3 of
\cref{met:b2:cpp-templates:concept-first} were skipped: the operations of the body were not all listed. The right
\omterm{def:b2:cpp-templates:concept}{concept} also requires \code{{(a + b) / 2} -> std::convertible_to<T>} in its \code{requires} expression; strings are
then refused at the call.
\end{solution}

\begin{solution}{exo:b2:cpp-templates:7}
\omcode{code/bachelor-2/04-cpp-templates/exercises/bsearch.hpp}{6}{29}{Binary search over any random-access range}
Invariant: every element before \code{first} is less than the key, and every element from \code{first + n} on is not;
the answer is therefore in $[\code{first}, \code{first} + n]$. Each step keeps it and at least halves $n$, so the loop
ends after at most $\lceil \log_2(n + 1) \rceil$ steps with $n = 0$, and \code{first} is the answer. The test
compares with \code{std::lower_bound} and \code{std::binary_search} on 200 random sorted vectors and every key from
$-1$ to 41.
\end{solution}

\begin{solution}{exo:b2:cpp-templates:8}
\code{Pow<20>::value} needs \code{Pow<19>::value}, which needs \code{Pow<18>}, and so on down to \code{Pow<0>}, the
specialisation written by hand: 20 classes instantiated, \code{Pow<1>} to \code{Pow<20>}. With \code{Pow<25>} as well,
only \code{Pow<21>} to \code{Pow<25>} are new, since each class is instantiated once
(\cref{prop:b2:cpp-templates:instantiations}): 25 in all. A \code{constexpr} function \code{power(3, n)} is one
function, which the compiler runs on the constant arguments; no class at all is generated, and compile-time loops
are no longer limited by the depth of recursive instantiation.
\end{solution}

\begin{solution}{exo:b2:cpp-templates:9}
\omcode{code/bachelor-2/04-cpp-templates/exercises/matrix.hpp}{6}{33}{Dimensions in the type}
The product is a \omterm{def:b2:cpp-templates:template}{template} in four parameters, \code{T}, \code{R}, \code{K} and \code{C}, and its two parameters
share \code{K}: a $2 \times 3$ by a $2 \times 3$ would need $K = 3$ and $K = 2$ at once, deduction fails, and there is
no other \code{operator*}. The test defines the \omterm{def:b2:cpp-templates:concept}{concept} \code{Multipliable<A, B>} (the product is valid) and asserts
it for $2 \times 3$ by $3 \times 4$ and refutes it for $2 \times 3$ by $2 \times 3$, at compile time; it also checks
$M M^{\mathsf{T}}$ for a $2 \times 3$ matrix.
\end{solution}

\begin{solution}{exo:b2:cpp-templates:10}
\omcodes{code/bachelor-2/04-cpp-templates/exercises/heap.hpp}{10-29,30-49}{\code{Heap<T, Less>}}
\code{push} adds at the end and sifts up while the parent is ``less''; \code{pop} moves the last element to the root
and sifts it down, choosing the larger child. Each instantiation knows the comparator's \emph{type}, hence the exact
function it calls (for \code{std::less<int>}, the built-in \code{<}), and can inline it. \code{qsort} receives a
pointer \emph{value}, known only when the program runs, and is compiled once in the C library: it must make an
indirect call each time.
\end{solution}

\begin{solution}{exo:b2:cpp-templates:11}
(a) At most $kmb$ bytes: one copy of each member per type. (b) Each of the $u$ object files contains up to $km$
functions, $ukm$ weak copies in all; the program, after linking, $km$. The compile time grows with $u$, the program
does not. (c) 30 symbols for 3 types, about 10 functions per type (\omterm{def:b2:cpp-classes-raii:constructor}{constructors}, \omterm{def:b2:cpp-classes-raii:constructor}{destructor}, \code{push_back},
\code{reserve}, \code{clear}, \code{operator[]}, \code{begin}, \code{end} and others). A \code{void *} version is
compiled once, about 10 functions for every type together, at the price of the type checks and the inlining.
\end{solution}

\begin{solution}{exo:b2:cpp-templates:12}
\omcode{code/bachelor-2/04-cpp-templates/exercises/binomial.hpp}{6}{21}{Pascal's triangle at compile time}
The largest entry of row $n$ is the middle one. $\binom{67}{33} \approx 1.42 \times 10^{19}$ is below $2^{64} \approx
1.84 \times 10^{19}$, and $\binom{68}{34} \approx 2.85 \times 10^{19}$ is above: 67 is the last row that fits entirely.
Pascal's rule only adds entries of the previous row, and each sum is an entry of the table, so no intermediate
result overflows up to $N = 67$. The test checks the rule for every entry and the row sum
$\sum_k \binom{20}{k} = 2^{20}$.
\end{solution}

\begin{solution}{pb:b2:cpp-templates:1}
\textbf{1.} Read both values as \code{double} and return \code{(x > y) - (x < y)}. \code{x - y} is converted to
\code{int}: a difference of 0.5 becomes 0, ``equal''. For \code{int}, \code{x - y} overflows when, say, $x$ is near
\code{INT_MAX} and $y$ negative, which is undefined behaviour and in practice gives the wrong sign.

\textbf{2.} Compare \code{due} with the same three-way formula, and if equal, compare \code{id}.

\textbf{3.} \code{qsort} swaps elements as blocks of bytes. A \code{std::string} may point into its own small buffer
or own a heap block; a byte copy breaks the first (the pointer still points into the old place) and duplicates the
owner of the second. Moving C++ objects that are not trivially copyable without their \omterm{def:b2:cpp-classes-raii:constructor}{constructors} is undefined
behaviour.

\textbf{4.} Only the types of the four arguments: any object pointer converts to \code{void *}. \code{qsort(d, n,
sizeof(int), compare_ints)} on an array of \code{double} compiles and sorts garbage.

\textbf{5.} About $n \log_2 n = 10^6 \times 20 = 2 \times 10^7$ calls.

\textbf{6.} Nothing: no algorithm, no complexity, no stability.

\textbf{7.} For $i$ from 1 to $n - 1$: \code{T x = std::move(a[i])}; shift the elements before it that are greater
than \code{x} (\code{less(x, a[j - 1])}) one place right, by move-assignment; put \code{x} in the hole.

\textbf{8.} A lambda \code{(double x, double y) -> x < y}; \code{std::less<>()}; a lambda that compares \code{due},
then \code{id}.

\textbf{9.} Three functions: \code{insertion_sort} for \code{double} with the first lambda's type, for
\code{std::string} with \code{std::less<void>}, and for \code{loan} with the second lambda's type (every lambda has a
type of its own).

\textbf{10.} The \omterm{def:b2:cpp-templates:template}{template} is compiled for \code{std::string} and moves elements with \code{std::string}'s \omterm{def:b2:cpp-copy-move:move}{move
constructor} and assignment, which keep its invariant.

\textbf{11.} \code{T tmp = std::move(a); a = std::move(b); b = std::move(tmp);}, or \code{std::swap(a, b)}: three moves
of a few pointers each for long strings, whatever their length.

\textbf{12.} Up to $n^2 / 2 = 5 \times 10^{11}$ comparisons against about $2 \times 10^7$ for \code{qsort}: the
algorithm matters far more than inlining.

\textbf{13.} A \code{requires} expression with parameters \code{L less}, \code{const T &a} and
\code{const T &b}, requiring \code{less(a, b)} to be valid and convertible to \code{bool}; the library's
\code{std::predicate<L, const T &, const T &>} says the same.

\textbf{14.} \code{std::movable<T>}.

\textbf{15.} No: the \omterm{def:b2:cpp-templates:concept}{concept} calls it with \code{const T &}, which does not bind to \code{T &}. That is right: a
comparison must not modify its arguments, and the sort relies on it.

\textbf{16.} \code{<=} is not a strict order: \code{x <= x} is true. Sorts assume irreflexivity; insertion sort then
loses stability, and \code{std::sort} given such a comparator may read past the end of the range. A \omterm{def:b2:cpp-templates:concept}{concept} checks
syntax, not meaning: \code{std::strict_weak_order} only documents the requirement.

\textbf{17.} Take two iterators \code{first} and \code{last} of a type \code{It} constrained by the
library's \omterm{def:b2:cpp-templates:concept}{concept} of random-access iterators, and keep \code{Less less}; \code{std::sortable<It, Less>} states both
requirements at once.

\textbf{18.} For each type, sort a copy of random input and check \code{std::is_sorted} with the same comparator and
\code{std::is_permutation} against the input.

\textbf{19.} In the printed run, for $n = 10^6$: \code{qsort} takes of the order of 90\,ns per element, about 4\,ns
per comparison (dividing by $2 \times 10^7$); \code{std::sort} about half of each.

\textbf{20.} Yes, between about 1.5 and 2: the difference is a cost per comparison, not a different growth. Both make
$O(n \log n)$ comparisons; each comparison costs more in \code{qsort}.

\textbf{21.} \code{qsort}: pass two addresses, make an indirect call, load the two \code{int}s through the pointers,
compute the three-way result with two comparisons, return, and test that result. \code{std::sort}: one comparison of
two registers, and a branch.

\textbf{22.} The night: tens of milliseconds saved, irrelevant. Ten thousand orders every millisecond: from the
printed run a sort of $10^4$ elements takes about half a millisecond with \code{qsort}; the difference is a large
share of the budget, and the rewrite is worth it.

\textbf{23.} At \code{-O0} nothing is inlined: \code{std::sort}'s comparisons become calls of the lambda's
\code{operator()}, and its iterator operations calls too. The ratio falls to about 1 or below.

\textbf{24.} \emph{The comparator speed-up of \code{std::sort} over \code{qsort}, measured on the laptop}: in the
printed run, a factor of about 1.5 to 2 for $10^4$ to $10^7$ integers, constant in $n$, because the comparison is
inlined into the instantiated sort instead of being called through a pointer.
\end{solution}
